\providecommand{\placeholder}[1]{\textbf{[TODO--AUTHOR: #1]}}

\section{Supplementary Method Details}
\label{supp:method-details}

This supplement specifies the implementation behind the method in the main
paper.  We use \(c\) for a claim, \(a_i\) for its \(i\)-th atomic
verification unit, and \(u_j\) for a candidate evidence passage.  Unless
otherwise stated, all normalizations are performed independently within a
claim's candidate set.  We write
\(\operatorname{clip}(x)=\min(1,\max(0,x))\).

\subsection{Canonical Claim Atomization}
\label{supp:atomization}
% Implementation source: src/fact_checking/selectors/claim_atomization.py

\paragraph{Prompt and output schema.}
Claim decomposition is performed once, before retrieval, using an
OpenAI-compatible chat-completions endpoint.  The logical model name in the
frozen configuration is \texttt{deepseek-v4-flash}; the exact provider
snapshot and invocation interval should be recorded as
\placeholder{atomizer provider snapshot/revision and UTC invocation date
window}.  The request uses temperature \(0\), top-\(p=1\), a maximum of 2,048
generated tokens, JSON-object response mode, and disabled model ``thinking.''
The prompt version is \texttt{claim\_atomization\_v0\_1}, and the parsed schema
version is \texttt{claim\_atoms\_v1}.  The exact system prompt is:

\begin{lstlisting}
You create canonical atomic verification units for fact-checking.
Use only the given claim. Do not use outside knowledge.
Return strictly valid JSON.
\end{lstlisting}

The exact user template is:

\begin{lstlisting}
Claim:
{claim}

Decompose the claim into 1 to 6 atomic verification units.

Rules:
- Each atom must be a complete proposition with clear truth conditions.
- Do not create standalone entity, date, number, or topic fragments.
- Keep dates, quantities, negation, modality, attribution, comparison targets,
  offices, locations, and scope inside the proposition.
- Split only when the claim contains multiple separately verifiable
  propositions.
- A single-sentence claim making one factual assertion should usually have one
  atom.
- Preserve the claim's meaning; do not add facts not present in the claim.
- The query_rendering is only a retrieval view of the same atom, not an
  independent question.

Return JSON exactly in this schema:
{
  "complexity": "simple|moderate|complex",
  "claim_atoms": [
    {
      "atom_id": "A1",
      "proposition": "...",
      "importance": 1.0,
      "type":
        "entity|quantity|date|comparison|causal|attribution|policy|other",
      "keywords": ["..."],
      "query_rendering": "...?"
    }
  ]
}
\end{lstlisting}

\paragraph{Validation and deterministic fallback.}
The parser first attempts to decode the complete response as JSON and then,
if necessary, extracts the first balanced JSON object.  It requires between
one and six nonempty proposition objects.  Invalid complexity values become
\texttt{moderate}; invalid atom types become \texttt{other}; importance is
clipped to \([0.05,1]\); and a missing terminal question mark is appended to
\texttt{query\_rendering}.  Exact duplicate propositions are removed after
lowercasing and whitespace normalization, and the surviving atoms are
renumbered \(A1,A2,\ldots\).  Transient API failures are retried up to five
times after the initial request, and a schema-related generation failure can
trigger at most two additional generations.

If no valid output is obtained, the system emits one atom whose proposition
is the complete claim, whose importance is \(1\), and whose type is
\texttt{other}.  Its retrieval rendering is:

\begin{lstlisting}
What evidence verifies whether this claim is true: {claim}?
\end{lstlisting}

\subsection{Sentence Segmentation and Claim-Aware Chunking}
\label{supp:chunking}
% Implementation sources:
% src/fact_checking/build/chunking.py
% src/fact_checking/utils/text.py

Raw reports are whitespace-normalized and split with the regular expression:
\begin{lstlisting}
(?<=[.!?])\s+(?=["'(\[]*[A-Z0-9])
\end{lstlisting}
Before splitting, the following abbreviations are protected:
\texttt{mr.}, \texttt{mrs.}, \texttt{ms.}, \texttt{dr.}, \texttt{prof.},
\texttt{sr.}, \texttt{jr.}, \texttt{st.}, \texttt{vs.}, \texttt{etc.},
\texttt{i.e.}, \texttt{e.g.}, \texttt{u.s.}, \texttt{u.k.}, \texttt{no.},
and \texttt{fig.}.  Chunk length is measured with
\texttt{[A-Za-z0-9\%\$]+} word tokens.

For sentence \(s_i\), claim relevance is
\begin{equation}
\begin{aligned}
r_i={}&0.70\,\operatorname{mm}(\operatorname{cos}(c,s_i))\\
&+0.20\,\operatorname{mm}(\operatorname{LexF1}(c,s_i))\\
&+0.10\,\operatorname{mm}(\operatorname{BM25L}(c,s_i)),
\end{aligned}
\label{eq:supp-sentence-relevance}
\end{equation}
where \(\operatorname{mm}\) is within-report min--max scaling and returns zero
for a constant vector.  The boundary score between adjacent sentences is
\begin{equation}
\begin{aligned}
b_i={}&w_{\mathrm{sem}}
\left(1-\operatorname{cos}(h_i,h_{i+1})\right)\\
&+w_{\mathrm{rel}}\lvert r_i-r_{i+1}\rvert .
\end{aligned}
\label{eq:supp-boundary}
\end{equation}
The standard configuration uses
\(w_{\mathrm{sem}}=0.75\), \(w_{\mathrm{rel}}=0.25\), and selects a positive
local maximum only when
\(b_i\geq \operatorname{mean}(b)+0.5\operatorname{std}(b)\).
If \(s_{i+1}\) starts with
\{\textit{he, she, they, it, this, that, these, those}\}, \(0.10\) is
subtracted from \(b_i\), with a floor at zero, to discourage a split before a
likely anaphoric continuation.

Chunks contain at most three sentences, at most 150 word tokens, and normally
at least 20 word tokens.  A one-sentence chunk is retained when its relevance
is at least \(0.55\).  An overlong span is recursively split at its highest
internal boundary.  An undersized span is merged with a feasible adjacent
span by, in order, greater embedding similarity, greater neighbor relevance,
and a right-neighbor tie break; a retained high-relevance singleton is exempt
from this merge.

The resolved RAWFC-tight configuration uses the same algorithm with
\(w_{\mathrm{sem}}=0.65\), \(w_{\mathrm{rel}}=0.35\), a local-peak
threshold of \(\operatorname{mean}(b)+0.35\operatorname{std}(b)\), and at
most two sentences per chunk.  Its remaining length, singleton, and
coreference settings are unchanged.

\subsection{Hybrid Claim- and Atom-Conditioned Retrieval}
\label{supp:retrieval}
% Implementation sources:
% src/fact_checking/retrieval/embedder.py
% src/fact_checking/retrieval/text_utils.py
% src/fact_checking/retrieval/mmr.py
% src/fact_checking/selectors/atom_conditioned_retrieval.py
% src/fact_checking/selectors/atom_retrieval_union.py

\paragraph{Dense representation.}
The BGE encoder~\citep{Xiao2023BGE} prepends
\begin{lstlisting}
Represent this sentence for searching relevant passages:
\end{lstlisting}
to queries, but not to passages.  It uses \texttt{AutoModel}, attention-masked
mean pooling of the final hidden states, and \(L_2\) normalization.  The
maximum encoded length is 256 tokens and the batch size is 64.
\placeholder{frozen embedding checkpoint revision and retrieval precision by
run}.

\paragraph{Lexical scores.}
Lexical tokens are lowercased matches of
\texttt{[A-Za-z0-9\_'-]+}.  We remove the following fixed stopword set:
\begin{lstlisting}
a an the and or to of on in for by with is are was were be been being
that this it as at from will would could should can may might do does
did have has had not but if than then into their there about literally
\end{lstlisting}
For query and passage content-token multisets \(Q\) and \(D\), the overlap is
\(o=\sum_{t\in Q\cap D}\min(\operatorname{tf}_Q(t),
\operatorname{tf}_D(t))\).  We use precision \(o/|D|\), recall \(o/|Q|\), and
their harmonic mean \(\operatorname{LexF1}\).

The implementation's third route is deliberately named
\emph{BM25-like} (\(\operatorname{BM25L}\)): unlike corpus-statistical
BM25~\citep{Robertson2009BM25}, it does not estimate document frequency
from a corpus.  For each distinct content token in the query,
\begin{equation}
\begin{aligned}
\iota_t(D)
&=\log\!\left(1+\frac{1}{1+f_t}\right)+0.5,\\
\operatorname{BM25L}(Q,D)
&=\sum_{\substack{t\in\operatorname{uniq}(Q)\\f_t>0}}
\iota_t(D)\,
\frac{f_t(k_1+1)}
{f_t+k_1(1-b+b|D|/\overline d)},\\
k_1&=1.2,\quad b=0.75,\quad \overline d=18,
\end{aligned}
\label{eq:supp-bm25-like}
\end{equation}
where \(f_t=\operatorname{tf}_D(t)\).  Thus the bracketed factor is a local,
frequency-based surrogate, not a corpus-statistical inverse document
frequency.

\paragraph{Route scoring and fusion.}
For query \(q\) and candidate \(u\), the route score is
\begin{equation}
\begin{aligned}
H(q,u)={}&0.70\,\operatorname{mm}(d(q,u))\\
&+0.20\,\operatorname{mm}(\operatorname{LexF1}(q,u))\\
&+0.10\,\operatorname{mm}(\operatorname{BM25L}(q,u)).
\end{aligned}
\label{eq:supp-hybrid-route}
\end{equation}
Min--max scaling returns an all-zero vector when the score range is below
\(10^{-8}\).  Each atom uses its \texttt{query\_rendering}, falling back to
its proposition, and retains its top \(k_a\) routes.  Atom routes are fused
by reciprocal-rank fusion (RRF)~\citep{Cormack2009RRF}:
\begin{equation}
\begin{aligned}
\operatorname{RRF}(u)
&=\sum_{a:\,u\in R_a}
\frac{w_a}{60+\operatorname{rank}_a(u)},\\
w_a&=1.
\end{aligned}
\label{eq:supp-rrf}
\end{equation}
Passages are deduplicated by lowercasing and collapsing whitespace.  Before
truncation to \(k_{\mathrm{pool}}\), the atom pool is ordered by decreasing
RRF, decreasing maximum atom-route hybrid score, decreasing number of
matched atoms, and increasing original candidate index.

We then take the deduplicated union of a claim-level route and the atom pool.
Its pre-MMR relevance is
\begin{equation}
\begin{aligned}
q(u)={}&\max\!\left(H(c,u),\max_a H(a,u)\right)\\
&+\operatorname{RRF}(u)+0.004\,n_{\mathrm{route}}(u).
\end{aligned}
\label{eq:supp-union-relevance}
\end{equation}
After within-claim min--max scaling of \(q\), maximal marginal relevance
(MMR)~\citep{Goldstein1998MMR} chooses the next candidate according to
\begin{equation}
\begin{aligned}
&\lambda q(u)
-(1-\lambda)\max_{v\in S}\operatorname{cos}(u,v),\\
&\lambda=0.70.
\end{aligned}
\label{eq:supp-mmr}
\end{equation}
The frozen per-dataset retrieval sizes are
\placeholder{table of claim-route size, \(k_a\), atom-pool size, and final
union-pool size for each dataset}.

\subsection{Evidence Map Annotation}
\label{supp:evidence-map}
% Implementation source:
% src/fact_checking/selectors/evidence_map_selector.py

\paragraph{Pair-level annotation prompt.}
The Evidence Mapper receives the claim, the fixed atom list, and the retrieved
passages; it cannot alter the atoms.  The exact model snapshot and processing
interval are the DeepSeek API model \texttt{deepseek-v4-flash}, with response
system fingerprint
\texttt{fp\_8b330d02d0\_prod0820\_fp8\_kvcache\_20260402}; recorded
invocations span 2026-06-21 17:07:40 to 2026-07-10 20:41:13 UTC.  Requests use
temperature \(0\), top-\(p=1\), at
most 8,192 generated tokens, JSON-object response mode, and disabled model
``thinking.''  The system prompt is:

\begin{lstlisting}
You are a fact-checking evidence analyst. Use only the claim, fixed atomic
verification units, and evidence passages supplied by the user. Do not
create, remove, merge, split, or rename atoms. Return valid JSON only.
\end{lstlisting}

The user prompt instantiates the following exact schema and field policy:

\begin{lstlisting}
Task: map each evidence passage to the fixed atomic verification units.
Do not create, remove, merge, split, or rename atoms.
Return one JSON object. Use this schema and do not add other fields:
{"candidate_atom_alignments":[{"evidence_id":"E01","atom_id":"A1",
"relation":"support|refute|qualify|mixed|insufficient|background|irrelevant",
"directness":"direct|partial|context|none",
"evidence_role":"primary_support|primary_refute|partial_support|
partial_refute|qualifying_context|background_context|duplicate|irrelevant",
"key_spans":["short quote"],"duplicate_group":"G1","confidence":0.0}]}

Field guide:
- candidate_atom_alignments.evidence_id: one of the evidence IDs below.
- candidate_atom_alignments.atom_id: one of the fixed atom IDs below.
- candidate_atom_alignments.relation: support, refute, qualify, mixed,
  insufficient, background, or irrelevant.
- candidate_atom_alignments.directness: direct, partial, context, or none.
- candidate_atom_alignments.evidence_role: primary_support, primary_refute,
  partial_support, partial_refute, qualifying_context, background_context,
  duplicate, or irrelevant.
- candidate_atom_alignments.key_spans: short substrings copied from the
  passage.
- candidate_atom_alignments.duplicate_group: same group ID for near-duplicate
  passages; use an empty string if none.
- candidate_atom_alignments.confidence: 0.0 to 1.0 for how certain the
  pair-level alignment is.

Output policy:
- Omit irrelevant evidence-atom pairs instead of listing them.
- If an evidence passage has no useful relation to any atom, output no row
  for that passage.
- Return at most one row for each evidence_id and atom_id pair.
- If no informative pairs exist, return
  {"candidate_atom_alignments":[]}.

Claim:
{claim}

Fixed atomic verification units:
A1: {atom proposition}
...

Evidence passages:
E01: {passage}
...
\end{lstlisting}

\paragraph{The complete annotation ontology.}
The mapper uses seven relation labels and four directness labels.  These are
the actual annotation categories; the four relation rows in the main paper
are an abstract transition representation obtained by grouping them.

\begin{table}[t]
\centering
\small
\begin{tabularx}{\linewidth}{@{}p{0.22\linewidth}X
p{0.24\linewidth}@{}}
\hline
Mapper relation & Pair-level meaning & Transition group \\
\hline
\texttt{support} & Evidence favors the atom. & support \\
\texttt{refute} & Evidence contradicts the atom. & refute \\
\texttt{qualify} & Evidence adds a condition or material qualification. &
qualify \\
\texttt{mixed} & The same passage contains materially supporting and
opposing/qualifying content. & qualify \\
\texttt{insufficient} & The passage concerns the atom but cannot resolve it. &
non-resolving \\
\texttt{background} & Context related to the atom without resolving it. &
non-resolving \\
\texttt{irrelevant} & No useful relation to the atom. & non-resolving \\
\hline
\end{tabularx}
\caption{Full Evidence Mapper relation ontology and its state-machine
projection.}
\label{tab:supp-map-relations}
\end{table}

\begin{table}[t]
\centering
\small
\begin{tabularx}{\linewidth}{lXc}
\hline
Directness & Meaning & Numeric factor \\
\hline
\texttt{direct} & Directly resolves the targeted proposition. & \(1.00\) \\
\texttt{partial} & Addresses only part of the proposition or resolves it
indirectly. & \(0.65\) \\
\texttt{context} & Supplies contextual material. & \(0.25\) \\
\texttt{none} & Supplies no resolving connection. & \(0.00\) \\
\hline
\end{tabularx}
\caption{Full directness ontology.}
\label{tab:supp-map-directness}
\end{table}

The role field takes one of
\texttt{primary\_support}, \texttt{primary\_refute},
\texttt{partial\_support}, \texttt{partial\_refute},
\texttt{qualifying\_context}, \texttt{background\_context},
\texttt{duplicate}, or \texttt{irrelevant}.

\paragraph{Validation, collapse, and fallback.}
Rows with unknown evidence or atom identifiers are discarded.  Invalid
relations and directness values become \texttt{irrelevant} and \texttt{none},
respectively; confidence is clipped to \([0,1]\); at most three copied spans
of at most 240 characters each are retained.  Malformed \texttt{key\_spans}
arrays receive a narrow JSON repair.  Duplicate pair rows are removed and
the result is sorted by evidence ID and atom ID.

For candidate-level features, pair rows are collapsed as follows.  Covered
atoms exclude only \texttt{background} and \texttt{irrelevant} pairs.  A
candidate containing both support and refute pairs is labeled
\texttt{mixed}; otherwise relation priority is
\[
\texttt{refute}>\texttt{qualify}>\texttt{mixed}>\texttt{support}>
\texttt{insufficient}>\texttt{background}>\texttt{irrelevant}.
\]
Directness priority is
\(\texttt{direct}>\texttt{partial}>\texttt{context}>\texttt{none}\).
Confidence is the maximum pair confidence, the first nonempty duplicate group
is used, and up to three distinct spans are retained.  A passage with no
valid pair receives the explicit fallback
\(\{\texttt{relation=irrelevant},\texttt{directness=none},
\texttt{confidence=0}\}\).  A missing or invalid map applies the same
fallback to every passage.

\paragraph{Candidate map quality.}
Let \(A(u)\) be the atoms covered by a candidate after collapse.  Atom
importance \(i_a\) is converted to
\[
\widehat i_a =
\begin{cases}
\operatorname{clip}_{[0.05,1]}(i_a/5),& i_a>1,\\
\operatorname{clip}_{[0.05,1]}(i_a),& i_a\leq 1,
\end{cases}
\]
and
\[
\operatorname{cov}(u)=
\frac{\sum_{a\in A(u)}\widehat i_a}{\sum_a\widehat i_a}.
\]
With directness score \(d\), polar-relation score \(p\), confidence \(\gamma\),
span-presence indicator \(s\), and background penalty \(g\), map quality is
\begin{equation}
\begin{aligned}
m(u)=\operatorname{clip}\!\bigl(&
0.35\operatorname{cov}(u)+0.25d+0.20p\\
&+0.10\gamma+0.10s-0.20g\bigr).
\end{aligned}
\label{eq:supp-map-quality}
\end{equation}
Here \(d\) is given in Table~\ref{tab:supp-map-directness};
\(p=1\) for support/refute, \(0.60\) for qualify/mixed, \(0.15\) for
background, and \(0\) otherwise.  The penalty is \(1\) for irrelevant,
\(0.9\) for background, \(0.8\) for directness none, \(0.55\) for context,
and \(0\) otherwise.

\subsection{Atom State and Transition Semantics}
\label{supp:state}
% Implementation sources:
% src/fact_checking/selectors/minimal_resolving_chain.py
% src/fact_checking/selectors/mrec_learned_marginal.py

Each atom has one hard state:
\[
\begin{aligned}
U&=\text{unresolved},\quad S=\text{supported},\\
R&=\text{refuted},\quad Q=\text{qualified},\\
C&=\text{conflicted}.
\end{aligned}
\]
For transition purposes, support maps to \(S\), refute to \(R\), and qualify
or mixed to \(Q\).  Insufficient, background, and irrelevant are
non-resolving.  The complete deterministic transition is:

\begin{table}[t]
\centering
\scriptsize
\setlength{\tabcolsep}{2.5pt}
\begin{tabular}{lccccc}
\hline
Incoming group & \(U\) & \(S\) & \(R\) & \(Q\) & \(C\) \\
\hline
support & \(S\) & \(S\) & \(C\) & \(Q\) & \(C\) \\
refute & \(R\) & \(C\) & \(R\) & \(Q\) & \(C\) \\
qualify/mixed & \(Q\) & \(Q\) & \(Q\) & \(Q\) & \(Q\) \\
non-resolving & \(U\) & \(S\) & \(R\) & \(Q\) & \(C\) \\
\hline
\end{tabular}
\caption{Atom-state transition function.}
\label{tab:supp-state-transition}
\end{table}

A passage may be aligned to multiple atoms, but a trace step has one target
pair and updates only that atom.  For a selected candidate, each valid pair is
assigned an operation and a target utility
\begin{equation}
v(a,u)=\pi_{\mathrm{op}}+10\,i_a+m(u)+0.1\,b(u),
\label{eq:supp-target-pair}
\end{equation}
where \(b(u)\) is the available base retrieval score and
\(\pi_{\mathrm{op}}\) is \(100,80,35,10,0\) for
\textsc{Open}, \textsc{Contrast}, \textsc{Corroborate}, \textsc{Bridge}, and
\textsc{Fallback}, respectively.  Ties are broken by relation priority
refute \(>\) qualify/mixed \(>\) support \(>\) insufficient, then map quality,
base score, and atom ID.  \textsc{Open} resolves an unresolved atom;
\textsc{Contrast} creates \(C\) for opposed support/refute or \(Q\) for a
qualification; \textsc{Corroborate} preserves a matching resolved state; and
\textsc{Bridge}/\textsc{Fallback} preserve the prior state.

\subsection{Twelve-Dimensional Marginal Utility}
\label{supp:marginal-features}
% Implementation source:
% src/fact_checking/selectors/mrec_learned_marginal.py

At selection step \(t\), every eligible candidate is represented by
\(\phi_t(u)\in[0,1]^{12}\).  Let \(p_U(a)\) be the current unresolved mass,
let
\[
\delta(r)\in\{1,0.65,0.25,0\}
\]
be the directness factor, and let \(\gamma_{ua}\) be pair confidence.  For
each resolving pair, define
\[
\begin{aligned}
g_{ua}&=\operatorname{clip}\!\left(
p_U(a)\,\delta(r_{ua})\,
\max(\gamma_{ua},0.5)\right),\\
g_a(u)&=\max_{(u,a)}g_{ua}.
\end{aligned}
\]
The features are:

\begin{enumerate}
\item \texttt{resolution\_delta}:
  \(\sum_a g_a(u)/|A|\).
\item \texttt{entropy\_reduction}: unresolved mass moved to a resolving
  state, averaged over atoms.  With the current one-hot state realization,
  this equals the preceding gain numerically.
\item \texttt{new\_atom\_coverage}: number of candidate-covered atoms not
  covered by an earlier selected step, divided by \(|A|\).
\item \texttt{new\_relation\_for\_atom}: one iff any candidate pair introduces
  a relation not previously observed for that atom.
\item \texttt{stance\_tension}: the maximum
  \(\delta(r)\max(\gamma,0.5)\) for a qualifying relation after any resolved
  state, support after refute/conflict, or refute after support/conflict.
\item \texttt{corroboration\_gain}: the maximum
  \(\delta(r)\max(\gamma,0.5)\) for a previously seen resolving relation,
  multiplied by \(\max(\text{source novelty},\text{text novelty})\).
\item \texttt{source\_novelty}: zero iff a nonempty source key has already
  appeared, and one otherwise.
\item \texttt{text\_novelty}: zero iff either normalized evidence text or its
  nonempty duplicate-group ID has already appeared, and one otherwise.
\item \texttt{map\_confidence}: maximum pair confidence for the candidate.
\item \texttt{map\_quality}: \(m(u)\) from
  Eq.~\eqref{eq:supp-map-quality}.
\item \texttt{retrieval\_score}: the first available value among hybrid,
  baseline-hybrid, and base score, clipped to \([0,1]\).
\item \texttt{cost\_ratio}: token cost divided by the configured token budget;
  when no budget is configured, the maximum candidate cost in the pool is
  the denominator.
\end{enumerate}

\subsection{Scorer Fitting and Greedy Selection}
\label{supp:scorer-selector}
% Implementation sources:
% src/fact_checking/selectors/mrec_learned_marginal.py
% src/fact_checking/selectors/minimal_resolving_chain.py

\paragraph{Pairwise scorer fitting.}
The marginal scorer is
\begin{equation}
\begin{aligned}
F_\theta(u\mid z_t)
={}&\beta+\sum_{k=1}^{11}w_k\phi_{t,k}(u)\\
&-w_c\phi_{t,12}(u),\qquad w_k,w_c\geq0.
\end{aligned}
\label{eq:supp-marginal-score}
\end{equation}
Nonnegativity is enforced with a softplus parameterization.  The bias is
fixed at zero.  At each training rollout step, the deterministic
structure-only ranking tuple is, in descending preference,
\[
\begin{aligned}
(&\text{direct resolving},\\
&\text{partial/non-direct resolving},\\
&\text{new-atom coverage},\ \text{new relation},\\
&\text{map quality},\ \text{retrieval score},\\
&\text{earlier original index}).
\end{aligned}
\]
The first candidate is paired against every remaining candidate.  For a
preferred/nonpreferred pair \((u^+,u^-)\), the RankNet-style full-batch
objective~\citep{Burges2005RankNet} is
\[
\begin{aligned}
\Delta_\theta(u^+,u^-)
&=F_\theta(u^+)-F_\theta(u^-),\\
\mathcal{L}
&=\frac{1}{|\mathcal{P}|}
\sum_{(u^+,u^-)\in\mathcal{P}}
\operatorname{softplus}\!\left(
-\Delta_\theta(u^+,u^-)\right).
\end{aligned}
\]
Training uses the top 20 candidates, five rollout steps, and full-batch CPU
Adam with learning rate \(0.05\) for 30 epochs in the reference fit.
All positive weights and the cost weight are initialized to one; no
regularizer, scheduler, or bias update is used.  After each rollout winner,
state is replayed with the candidate's pair having highest directness, then
highest confidence, then lexicographically smallest atom ID.

\paragraph{Stateful selector pseudocode.}
The deployed ranking recomputes features after every selection:

\begin{enumerate}
\item Initialize all atom states to \(U\); initialize the selected set, source
  set, normalized-text set, duplicate-group set, and observed
  atom--relation set to empty.
\item Form the eligible set by excluding already selected candidates and any
  candidate whose normalized text or nonempty duplicate group has already
  been selected.  In budgeted modes, also exclude a candidate that would
  exceed the remaining token budget.
\item For every eligible \(u\), compute \(\phi_t(u)\) from the current state
  and score it with Eq.~\eqref{eq:supp-marginal-score}.
\item Choose the maximum-score candidate.  Deterministic tie breaks use, in
  order: larger \texttt{resolution\_delta}, larger
  \texttt{entropy\_reduction}, larger \texttt{stance\_tension}, larger
  \texttt{new\_atom\_coverage}, larger
  \texttt{new\_relation\_for\_atom}, larger
  \texttt{map\_confidence}, lower token cost, and lower original index.
\item Select the candidate's target evidence--atom pair using
  Eq.~\eqref{eq:supp-target-pair}; append the trace record; update only the
  target atom's hard and soft states; and update the novelty and observed
  relation sets.
\item Repeat until the candidate pool is exhausted, or until an explicitly
  configured trace horizon, utility threshold, or token budget is reached.
\end{enumerate}

\paragraph{Why a full ordering is retained.}
For the reported full-ordering configuration, the trace horizon is the
candidate-pool size.  The resulting \(T_{\mathrm{full}}\) is an internal audit
artifact, not the evidence sequence automatically exposed downstream.
Retaining it permits controlled capacity ablations and sampling of multiple
prefix lengths from exactly the same stateful ordering, without rerunning the
selector and thereby changing earlier decisions.  Operationally, the
implementation therefore builds the complete ordering first and truncates it
only when materializing an ablation or verifier prompt.

\subsection{Trace Projection and Prompt Rendering}
\label{supp:projection}
% Implementation sources:
% src/fact_checking/build/prompts.py
% scripts/phase5_selectors/build/build_trace_verifier_data.py

The capacity controller projects \(T_{\mathrm{full}}\) to a prefix.  Each
trace step stores \texttt{target\_resolved}, which is set when that step's
replayed resolved-atom rate reaches the configured
\(\rho_{\mathrm{target}}\).  Under the min--max rule, the renderer scans in
order and, after at least \(K_{\min}\) steps, stops at the first step carrying
this flag, subject to \(K_{\max}\):
\[
\begin{aligned}
H&=\min\{|T_{\mathrm{full}}|,K_{\max}\},\\
\mathcal R&=\{t:K_{\min}\leq t\leq H,\
\rho_t\geq\rho_{\mathrm{target}}\}.
\end{aligned}
\]
\begin{equation}
K_{\mathrm{count}}=
\begin{cases}
\min\mathcal R,&\mathcal R\neq\varnothing,\\
H,&\mathcal R=\varnothing,
\end{cases}
\label{eq:supp-count-projection}
\end{equation}
Fixed-count and token-cap controllers use the corresponding longest feasible
prefix.  The
reference count policies are listed in
Table~\ref{tab:supp-lora-settings}; the complete RQ3 controller grid is
listed in Table~\ref{tab:supp-capacity-grid}.

Prompt construction is tokenizer-aware.  The maximum training sequence length
is 1,024 tokens; after reserving target tokens, complete evidence units are
removed from the tail until the chat prompt fits.  Thus ordinary truncation
preserves trace order and never removes an interior unit.  If the fixed prompt
plus the sole remaining evidence unit is still too long, an emergency
fallback retains the longest tokenizer prefix of that evidence that fits;
if even an empty evidence item cannot fit, no evidence is retained.

For the minimal trace view, selected step \(t\) is rendered exactly as
\begin{lstlisting}
Check: <step-specific cue>
<raw evidence text>
\end{lstlisting}
in selector order.  The renderer uses the stored cue and falls back to the
target atom proposition, then to a generic verification cue, when it is
missing.  The same atom cue can therefore occur in multiple steps, but every
evidence unit has exactly one cue.  Operation, state, transition reason,
confidence, relation, directness, and other trace metadata remain in the
audit artifact and are not inserted into this prompt view.

\FloatBarrier
